\documentclass[11pt,a4paper]{article}

% ─── Paquetes ────────────────────────────────────────────────────────────────
\usepackage[utf8]{inputenc}
\usepackage[T1]{fontenc}
% es-noshorthands desactiva los caracteres activos ("., <, >) que babel-spanish
% activa por defecto y que rompen el parser de pgfkeys/TikZ.
\usepackage[spanish,es-noshorthands]{babel}
\usepackage[margin=2.3cm]{geometry}
\usepackage{xcolor}
\usepackage{listings}
\usepackage{enumitem}
\usepackage{hyperref}
\usepackage{tabularx}
\usepackage{booktabs}
\usepackage{titlesec}
\usepackage{fancyhdr}
\usepackage{graphicx}

% ─── Estilo general ──────────────────────────────────────────────────────────
\hypersetup{
  colorlinks=true,
  linkcolor=blue!60!black,
  urlcolor=blue!60!black,
  citecolor=blue!60!black,
  plainpages=false,
  pdfpagelabels
}

% Evita el warning de fancyhdr "headheight is too small".
\setlength{\headheight}{14pt}
\setlength{\topmargin}{-3pt}

\pagestyle{fancy}
\fancyhf{}
\fancyhead[L]{\textbf{LudoElixir}}
\fancyhead[R]{Documentación Técnica}
\fancyfoot[C]{\thepage}

\titleformat{\section}{\Large\bfseries\color{blue!50!black}}{\thesection}{1em}{}
\titleformat{\subsection}{\large\bfseries\color{blue!40!black}}{\thesubsection}{1em}{}

% ─── Listado de código Elixir ────────────────────────────────────────────────
\definecolor{elxbg}{HTML}{F7F7FA}
\definecolor{elxkw}{HTML}{8E44AD}
\definecolor{elxstr}{HTML}{27AE60}
\definecolor{elxcom}{HTML}{7F8C8D}
\definecolor{elxatom}{HTML}{C0392B}

\lstdefinelanguage{elixir}{
  morekeywords={def,defp,defmodule,do,end,if,else,case,cond,when,
                use,alias,import,require,fn,true,false,nil,
                with,for,in,not,and,or,receive,send,spawn,
                @impl,@doc,@moduledoc},
  sensitive=true,
  morecomment=[l]{\#},
  morestring=[b]",
}

\lstset{
  language=elixir,
  basicstyle=\ttfamily\small,
  keywordstyle=\color{elxkw}\bfseries,
  commentstyle=\color{elxcom}\itshape,
  stringstyle=\color{elxstr},
  backgroundcolor=\color{elxbg},
  frame=single,
  rulecolor=\color{gray!30},
  framesep=6pt,
  breaklines=true,
  showstringspaces=false,
  columns=fullflexible,
  tabsize=2,
  numbers=left,
  numberstyle=\tiny\color{gray},
  xleftmargin=2em,
  literate=
    {á}{{\'a}}1 {é}{{\'e}}1 {í}{{\'\i}}1 {ó}{{\'o}}1 {ú}{{\'u}}1
    {Á}{{\'A}}1 {É}{{\'E}}1 {Í}{{\'I}}1 {Ó}{{\'O}}1 {Ú}{{\'U}}1
    {ñ}{{\~n}}1 {Ñ}{{\~N}}1 {¿}{{?`}}1 {¡}{{!`}}1
}

% ─── Documento ───────────────────────────────────────────────────────────────
\begin{document}

\begin{titlepage}
  \centering
  \vspace*{3cm}
  {\Huge\bfseries LudoElixir\par}
  \vspace{0.6cm}
  {\Large Documentación Técnica\par}
  \vspace{0.4cm}
  {\large Juego de Ludo multijugador en tiempo real\par}
  \vspace{0.2cm}
  {\large Elixir~/~Phoenix LiveView~/~OTP\par}
  \vspace{0.6cm}
  {\large
  Víctor Ramos \\[2pt]
  Albert Cáceres \\[2pt]
  Daniel Cayo \\[2pt]
  Juan Callo
  \par}
  \vfill
  \begin{tabular}{rl}
    \textbf{Lenguaje:}     & Elixir \texttt{\textasciitilde{}> 1.15} (OTP 26+) \\
    \textbf{Framework web:} & Phoenix \texttt{\textasciitilde{}> 1.8} + LiveView \\
    \textbf{Concurrencia:}  & Modelo de actores (GenServer, Registry, PubSub) \\
    \textbf{Repositorio:}   & \url{https://github.com/vjrm123/LudoElixir.git} \\
  \end{tabular}
  \vfill
  {\large \today\par}
\end{titlepage}

\tableofcontents
\newpage

% ─────────────────────────────────────────────────────────────────────────────
\section{Diagrama de la Estructura del Juego}

El sistema se organiza en tres capas claramente separadas: un \textbf{árbol de
supervisión OTP} que gestiona el ciclo de vida de los procesos, una
\textbf{capa de lógica pura} sin efectos secundarios que define las reglas
del juego, y una \textbf{capa de presentación} basada en Phoenix LiveView
que renderiza el estado en el navegador del jugador mediante WebSockets.

\subsection{Árbol de Supervisión}

El módulo \texttt{Ludo.Application} arranca el supervisor raíz
\texttt{Ludo.Supervisor} con estrategia \texttt{:one\_for\_one}. Cada hijo
es responsable de un aspecto independiente del sistema: telemetría,
descubrimiento de nodos (DNSCluster), bus de mensajes (PubSub), un registro
de salas, un supervisor dinámico para los procesos de partida y el endpoint
HTTP/WS de Phoenix.

\begin{figure}[htbp]
\centering
\includegraphics[width=\textwidth]{supervision_tree.png}
\caption{Árbol de supervisión OTP.}
\label{fig:supervision}
\end{figure}

\noindent
\textbf{Estrategia \texttt{:one\_for\_one}}: si un \texttt{GameServer} cae
(por ejemplo, por un error inesperado), solo ese proceso es reiniciado por
el \texttt{DynamicSupervisor}; las demás salas siguen jugando sin
interrupción. El estado de la sala caída se pierde (no se persiste a
disco), pero los demás jugadores no son afectados.

\subsection{Capa de Lógica Pura (\texttt{lib/ludo/})}

Esta capa contiene todo el conocimiento de las reglas del Ludo sin
acoplarse a procesos, mensajería ni a la web. Cada módulo expone funciones
puras (idempotentes y sin efectos), lo que hace que sea trivial de testear
y razonar.

\begin{center}
\renewcommand{\arraystretch}{1.3}
\begin{tabularx}{\linewidth}{l X}
\toprule
\textbf{Módulo} & \textbf{Responsabilidad} \\
\midrule
\texttt{Ludo.Estado}   & Struct con \texttt{codigo, host\_id, jugadores, modo, fase, turno\_idx, dado, tablero, fichas\_movibles}. Helpers de avance de turno y finalización. \\
\texttt{Ludo.Reglas}   & Motor del juego: \texttt{fichas\_movibles/4}, \texttt{aplicar\_movimiento/5}, \texttt{pasos\_de\_movimiento/3}, equipos y victoria. \\
\texttt{Ludo.Board}    & Geometría del tablero: 52 celdas del camino, 5 celdas por pasillo de llegada, 8 celdas seguras, casillas de salida por color. \\
\texttt{Ludo.Colores}  & Lista canónica de colores (\texttt{:rojo, :azul, :verde, :amarillo}) y mapa hexadecimal. \\
\texttt{Ludo.Salas}    & Facade pública: crear, unirse, salir, iniciar partida; valida entradas y genera códigos únicos de sala. \\
\bottomrule
\end{tabularx}
\end{center}

\subsection{Capa de Presentación (\texttt{lib/ludo\_web/})}

Tres LiveViews cubren todo el ciclo de juego, comunicadas por WebSocket con
el GenServer de su sala mediante \texttt{Phoenix.PubSub}.

\begin{center}
\renewcommand{\arraystretch}{1.3}
\begin{tabularx}{\linewidth}{l l X}
\toprule
\textbf{Ruta} & \textbf{LiveView} & \textbf{Función} \\
\midrule
\texttt{GET /}                  & \texttt{InicioLive}   & Landing: crear sala o unirse con código. \\
\texttt{GET /lobby/:codigo}     & \texttt{LobbyLive}    & Sala de espera; el host inicia la partida. \\
\texttt{GET /tablero/:codigo}   & \texttt{TableroLive}  & Tablero 15$\times$15 con fichas, dado y turnos. \\
\bottomrule
\end{tabularx}
\end{center}

\noindent
Tres \textit{hooks} JavaScript (\texttt{assets/js/app.js}) refinan la
experiencia de cliente:

\begin{itemize}[leftmargin=*, itemsep=2pt]
  \item \texttt{JugadorSession}: persiste el \texttt{jugador\_id} en
    \texttt{sessionStorage} para sobrevivir a recargas.
  \item \texttt{BoardHook}: anima el desplazamiento de fichas paso a paso
    mediante un \textit{ghost token}.
  \item \texttt{DiceHook}: anima el dado con rotaciones 3D antes de
    revelar el resultado real.
\end{itemize}

\subsection{Flujo de Datos}

\begin{figure}[htbp]
\centering
\includegraphics[width=0.85\textwidth]{dataflow.png}
\caption{Flujo de datos: desde el clic del usuario hasta la animación en el navegador.}
\label{fig:dataflow}
\end{figure}

\newpage
% ─────────────────────────────────────────────────────────────────────────────
\section{Estrategia de Concurrencia}

La concurrencia se apoya en tres pilares de OTP: un \textbf{GenServer}
por sala (serializa mensajes vía buzón FIFO), \textbf{Phoenix.PubSub}
para difundir cambios a todos los jugadores, y \textbf{temporizadores}
(\texttt{Process.send\_after/3}) para timeouts. La siguiente figura
muestra la interacción entre dos jugadores y el GenServer:

\begin{figure}[htbp]
\centering
\includegraphics[width=0.6\linewidth]{diagrama_secuencia.png}
\caption{Diagrama de secuencia: llamadas concurrentes al GenServer.}
\label{fig:concurrencia}
\end{figure}

\newpage
\vspace{-2mm}
\subsection{Funcionamiento}

\begin{enumerate}[leftmargin=*, itemsep=2pt]
  \item \textbf{Suscripción}: al montarse, cada LiveView (jugador) se
    suscribe al topic \texttt{"sala:CODIGO"} de \texttt{Phoenix.PubSub}.
  \item \textbf{Tirar dado}: el jugador activo envía un
    \texttt{GenServer.call(:tirar\_dado)}. El GenServer valida el turno,
    genera 1--6 y calcula las fichas movibles. Luego responde al
    llamador con \texttt{\{:ok, resultado\}} y \texttt{broadcast}ea el
    nuevo estado a todos los suscriptores.
  \item \textbf{Mover ficha}: similar: el jugador envía
    \texttt{call(:mover\_ficha)}, el GenServer aplica
    \texttt{Reglas.aplicar\_movimiento/5}, responde y difunde.
  \item \textbf{Intento simultáneo}: si dos jugadores llaman al mismo
    tiempo, el buzón FIFO del GenServer serializa las peticiones. El
    segundo llamador recibe \texttt{\{:error, :no\_es\_tu\_turno\}} porque
    el turno avanzó tras procesar la primera.
\end{enumerate}

\subsection{Resumen}

\begin{itemize}[leftmargin=*, itemsep=2pt]
  \item \textbf{GenServer.call/3} (síncrono): \texttt{tirar\_dado},
    \texttt{mover\_ficha}, \texttt{iniciar}, \texttt{unirse}.
  \item \textbf{GenServer.cast/2} (asíncrono): \texttt{salir}.
  \item \textbf{Process.send\_after/3} (auto-mensaje): timeout de turno
    (30~s) y auto-pase (2~s).
  \item \textbf{Phoenix.PubSub}: difusión 1 a N del nuevo estado a todas
    las LiveViews en paralelo.
  \item \textbf{Sin locks}: el buzón FIFO del GenServer garantiza
    exclusión mutua; las funciones de \texttt{Reglas} son puras y no
    comparten estado.
\end{itemize}

\section{Definición de las Principales Funciones}

\subsection{\texttt{Ludo.Reglas} --- motor del juego (funciones puras)}

Funciones auxiliares invocadas desde \texttt{aplicar\_movimiento/5}:

\begin{center}
\renewcommand{\arraystretch}{1.3}
\begin{tabularx}{\linewidth}{l X}
\toprule
\textbf{Función} & \textbf{Descripción} \\
\midrule
\texttt{get\_color/2}        & (privada) Obtiene el \texttt{color} de un jugador a partir de su \texttt{jugador\_id}. \\
\texttt{nueva\_posicion/3}   & (privada) Calcula la coordenada de destino tras avanzar \texttt{dado} casillas. \\
\texttt{resolver\_capturas/4} & (privada) Procesa capturas enemigas en la celda de destino, respetando celdas seguras y bloqueos. \\
\texttt{todas\_en\_meta?/2}  & (privada) Verifica que las 4 fichas del jugador estén en la meta. \\
\texttt{equipo\_gana?/3}     & Comprueba si el equipo completo del jugador ha ganado (modo \texttt{:equipos}). \\
\texttt{maybe\_add/3}        & (privada) Agrega un evento a la lista solo si la condición es \texttt{true}. \\
\bottomrule
\end{tabularx}
\end{center}

\begin{lstlisting}[caption={Esqueleto de \texttt{aplicar\_movimiento/5}}]
def aplicar_movimiento(tablero, jugador_id, ficha_id, dado, jugadores) do
  fichas = Map.get(tablero, jugador_id, [])

  case Enum.find(fichas, &(&1.id == ficha_id)) do
    nil -> {:error, :ficha_no_encontrada}

    ficha ->
      color      = get_color(jugadores, jugador_id)
      nueva_pos  = nueva_posicion(ficha.pos, color, dado)
      tablero2   = actualizar_ficha(tablero, jugador_id, ficha_id, nueva_pos)
      {tab3, caps} = resolver_capturas(tablero2, jugador_id, nueva_pos, jugadores)

      eventos =
        []
        |> maybe_add(:ficha_capturada, caps != [])
        |> maybe_add(:ficha_en_meta,   nueva_pos == :meta)
        |> maybe_add(:jugador_gana,    todas_en_meta?(tab3, jugador_id))
        |> maybe_add(:equipo_gana,     equipo_gana?(tab3, jugadores, jugador_id))

      {:ok, tab3, eventos}
  end
end
\end{lstlisting}

\subsection{\texttt{Ludo.GameServer} --- proceso por sala}

Funciones auxiliares invocadas desde \texttt{tirar\_dado}:

\begin{center}
\renewcommand{\arraystretch}{1.3}
\begin{tabularx}{\linewidth}{l l X}
\toprule
\textbf{Función} & \textbf{Tipo} & \textbf{Descripción} \\
\midrule
\texttt{tirar\_dado/2}       & call    & API pública que delega al GenServer mediante \texttt{GenServer.call/3}. \\
\texttt{Reglas.fichas\_movibles/4} & pública & Calcula las fichas que el jugador activo puede mover con el resultado del dado. \\
\texttt{broadcast!/2}        & (privada) & Difunde un mensaje a todos los suscriptores vía \texttt{Phoenix.PubSub}. \\
\bottomrule
\end{tabularx}
\end{center}

\begin{lstlisting}[caption={Manejo de \texttt{:tirar\_dado} con validaciones}]
def handle_call({:tirar_dado, jugador_id}, _from, estado) do
  jugador_en_turno = Enum.at(estado.jugadores, estado.turno_idx)

  cond do
    estado.fase != :jugando       -> {:reply, {:error, :partida_no_activa},   estado, @timeout_sala}
    jugador_en_turno.id != jugador_id -> {:reply, {:error, :no_es_tu_turno},  estado, @timeout_sala}
    estado.dado != nil            -> {:reply, {:error, :ya_tiraste},          estado, @timeout_sala}

    true ->
      resultado = Enum.random(1..6)
      movibles  = Reglas.fichas_movibles(estado.tablero, jugador_id,
                                          jugador_en_turno.color, resultado)
      nuevo = %{estado | dado: resultado, fichas_movibles: movibles}

      broadcast!(nuevo.codigo, {:dado_tirado, resultado, jugador_id, nuevo})
      {:reply, {:ok, resultado}, nuevo, @timeout_sala}
  end
end
\end{lstlisting}

\subsection{\texttt{Ludo.Estado} --- struct y helpers de turno}

\begin{lstlisting}[caption={Helpers de avance de turno}]
def avanzar_turno(%__MODULE__{} = estado, dado) do
  n = length(estado.jugadores)

  if dado == 6 and estado.fase == :jugando do
    %{estado | dado: nil, fichas_movibles: []}
  else
    %{estado | dado: nil, fichas_movibles: [],
               turno_idx: rem(estado.turno_idx + 1, n)}
  end
end
\end{lstlisting}

\subsection{\texttt{Ludo.Board} --- geometría del tablero}

Contiene las constantes del camino de 52 celdas, las coordenadas de cada celda (\texttt{cell\_coords/1}), los pasillos de llegada (\texttt{home\_lane\_coords/1}), las casillas de salida por color (\texttt{casilla\_salida/1}) y el conjunto de 8 celdas seguras (\texttt{celda\_segura?/1}).

\subsection{\texttt{Ludo.Salas} --- API \textit{facade}}

Funciones auxiliares invocadas desde \texttt{crear\_sala/3}:

\begin{center}
\renewcommand{\arraystretch}{1.3}
\begin{tabularx}{\linewidth}{l X}
\toprule
\textbf{Función} & \textbf{Descripción} \\
\midrule
\texttt{validar\_modo/1}    & (privada) Rechaza modos distintos a \texttt{:clasico} o \texttt{:equipos}. \\
\texttt{validar\_color/1}   & (privada) Rechaza colores fuera de \{\texttt{:rojo}, \texttt{:azul}, \texttt{:verde}, \texttt{:amarillo}\}. \\
\texttt{validar\_nombre/1}  & (privada) Rechaza nombres vacíos o menores a 2 caracteres. \\
\texttt{codigo\_unico/0}    & (privada) Genera un código de 6 caracteres no existente en el \texttt{Registry}. \\
\texttt{generar\_id/0}      & (privada) Genera un identificador único hexadecimal de 16 caracteres. \\
\texttt{arrancar\_game\_server/3} & (privada) Inicia un \texttt{GameServer} hijo del \texttt{DynamicSupervisor}. \\
\texttt{GameServer.unirse/2} & call & Agrega el jugador a la sala en el GenServer. \\
\bottomrule
\end{tabularx}
\end{center}

\begin{lstlisting}[caption={Pipeline de creación de sala con \texttt{with/1}}]
def crear_sala(nombre_jugador, color, modo \\ :clasico) do
  with :ok            <- validar_modo(modo),
       :ok            <- validar_color(color),
       :ok            <- validar_nombre(nombre_jugador),
       codigo         <- codigo_unico(),
       host_id        <- generar_id(),
       jugador        <- %{id: host_id, nombre: nombre_jugador, color: color},
       {:ok, _pid}    <- arrancar_game_server(codigo, host_id, modo),
       {:ok, estado}  <- GameServer.unirse(codigo, jugador) do
    {:ok, %{codigo: codigo, host_id: host_id, estado: estado}}
  end
end
\end{lstlisting}

\end{document}
